% TOP_PROBLEM: {"id": 5300020, "title": "Thurston's algorithm without critical finiteness", "category_id": 11, "category": {"id": 11, "name": "dynamical_systems", "display_name": "Dynamical Systems", "description": "Problems about long-term behavior of deterministic systems, Hamiltonian dynamics, and periodic orbits.", "slug": "dynamical-systems", "order_index": 11, "created_at": "2026-07-31T15:26:25.671Z"}, "status": "partially_solved"}
% FIRST_POSTED: null
% ENTRY_KIND: statement_only
% Imported from: batch04.tex; UnsolvedMath ID 5300020
\documentclass[11pt]{article}
\usepackage{fontspec}
\setmainfont{Latin Modern Roman}
\setsansfont{Latin Modern Sans}
\usepackage{amsmath,amssymb,mathtools}
\usepackage{unicode-math}
\setmathfont{Latin Modern Math}
\usepackage{xurl}
\usepackage[hidelinks]{hyperref}
\newcommand{\sref}[2]{\hyperlink{source:#1:#2}{[S#2]}}
\newcommand{\cataloguescope}[1]{\par\noindent\textbf{Question.} #1\par}
\begin{document}
% UnsolvedMath ID 5300020; source catalogue code AMR-052-0020
\setcounter{section}{5300019}
\section{Thurston's algorithm without critical finiteness}\label{5300020}
{\small\sffamily UnsolvedMath ID 5300020 · Dynamical Systems}
\subsection{Definitions and mathematical statement}

\paragraph{Shared notation.}$\mathbb N=\{1,2,\ldots\}$, and $\mathbb Z,\mathbb Q,\mathbb R,\mathbb C$ denote the integers, rationals, reals and complex numbers. Any other conventions are specified locally.


Identify $S^2$ with the Riemann sphere and choose three distinct base points. Starting from an orientation-preserving branched covering $f_0:S^2\to S^2$, let $h_j$ be the homeomorphism fixing those points that makes $r_j=f_j\circ h_j$ holomorphic, and put $f_{j+1}=h_j^{-1}\circ f_j\circ h_j$. Define $\phi_j=(h_0\circ\cdots\circ h_j)^{-1}$. A branched covering is locally a finite-degree covering except at finitely many branch points. No finiteness of the forward critical orbits is assumed. Uniform convergence on the sphere uses the spherical metric.
There is also an interval version. Let $I=[0,1]$ and let $f_0:I\to I$ be piecewise monotone with $d$ alternately increasing and decreasing monotonicity intervals, sending its endpoints to endpoints. At each step choose the orientation-preserving homeomorphism $h_j$ for which $p_j=f_j\circ h_j$ is a polynomial of degree $d$, and again set $f_{j+1}=h_j^{-1}\circ f_j\circ h_j$.
\paragraph{Statement recorded in the supplied source.}
Under what conditions do the rational maps $r_j$ converge uniformly to a rational map $r_\infty$? Under what conditions, and on what subset of $S^2$, do the coordinate maps $\phi_j$ converge uniformly? For the interval construction, likewise decide when and where the procedure converges. \sref{5300020}{2}
\subsection{Short English statement}
Determine conditions for convergence of the rational maps and coordinate changes produced by Thurston’s algorithm without assuming finite critical orbits, and determine convergence in its interval analogue.
\subsection{Sources}
\begin{description}
\item[\hypertarget{source:5300020:1}{S1}] ULAM.AI and the UnsolvedMath Contributors, UnsolvedMath: A Curated Collection of Open Mathematics Problems, record 5300020 (AMR-052-0020). CC BY 4.0. \url{https://huggingface.co/datasets/ulamai/UnsolvedMath}.
\item[\hypertarget{source:5300020:2}{S2}] John Milnor, Thurston’s algorithm without critical finiteness, in Ben Bielefeld and Mikhail Lyubich (editors), Problems in holomorphic dynamics (1992). Construction, displayed Problem, and Maps of the interval, pp. 12–13. \url{https://arxiv.org/pdf/math/9205209}.
\end{description}
\label{end:5300020}
\end{document}
